\section{Introduction}
\label{sec:introduction}

A family of Hamiltonian cycles can be viewed as a binary measurement system on
the edges of a graph: an edge receives one bit for every cycle, indicating
whether the cycle contains it.  Requiring every ordered pair \((e,f)\) to be
separated asks for a coordinate at which the signature of \(e\) is one and that
of \(f\) is zero.  This requirement is stronger than mere distinctness of the
signatures and leads naturally to an antichain condition.

This manuscript develops that question for Barnette graphs, namely simple,
planar, bipartite, cubic, and 3-connected graphs.  Its first purpose is to state
precisely what has been established by a finite computer-assisted proof.  Its
second purpose is to isolate a recurrent double-ladder structure among the
certified extremal examples.  The manuscript deliberately distinguishes four
levels of evidence:
\begin{enumerate}[label=(\roman*)]
  \item ordinary mathematical propositions proved in the text;
  \item finite computational results supported by explicit, independently
        checkable certificates;
  \item prospective computations whose values were frozen before evaluation;
        and
  \item open extremal conjectures beyond the certified census range.
\end{enumerate}

The certified computation gives all extremal values from order 8 through order
36 whenever the Barnette census is nonempty.  The order-10 census is empty and
odd orders are impossible for cubic graphs.  At the nonempty orders, exact
upper certificates are Hamiltonian edge-separating families, while lower bounds
are supplied either by packing requirements or by exhaustive exclusion over a
certified complete Hamiltonian-cycle universe.  Importantly, no greedy cover
size is used as an exact value.

The structurally recurring examples are denoted by \(D(a,b)\).  Two ladder
strips are joined by four fixed terminal edges.  The certified unique
maximizers at orders (8,12,16,20,24,28,32,36) are respectively
\[
 D(1,1),D(3,1),D(3,3),D(5,3),D(5,5),D(7,5),D(7,7),D(9,7).
\]
The tied extremal graphs at orders (22,26,30,34) have a different face
structure and are not forced into this uniform family.

The finite pattern leads to an infinite-family theorem.  A rail/rung
recurrence classifies all Hamiltonian cycles of \(D(a,b)\) for odd
\(a,b\ge3\).  A symbolic packing supplies a lower bound for \(\hsep\), while a
boundary-plus-parity family of Hamiltonian cycles supplies a matching upper
bound through an incidence-signature antichain.  Consequently,
\[
 \hsep(D(a,b))=\frac{(a+2)(b+2)-1}{2}
 \qquad(a,b\ge3\text{ odd}).
\]
This settles the invariant on the double-ladder family, not its extremality
among all Barnette graphs of the same order.

\paragraph{Literature-review status.}
\textbf{TODO:} A systematic peer-reviewed literature review is required before
any novelty claim is made.  It must cover separating and completely separating
systems, edge separation by cycles, Hamiltonian-cycle covers, Barnette graph
generation and census results, and reducible configurations or expansion
operations.  The notation \(\hsep\) and the descriptive name used here are
therefore provisional.

The remainder of the paper defines the invariant and its certificate logic,
describes the census and verification pipeline, states the certified extremal
sequence, analyzes the double-ladder family, proves matching symbolic bounds,
and separates prospective confirmations from open extremal conjectures.
